\documentclass[12pt,a4paper]{report}
\usepackage[utf8]{inputenc}
\usepackage[T1]{fontenc}
\usepackage[french]{babel}
\usepackage{geometry}
\geometry{top=2.5cm,bottom=2.5cm,left=2.5cm,right=2.5cm}
\usepackage{graphicx}
\usepackage{xcolor}
\usepackage{tikz}
\usetikzlibrary{shapes.geometric,arrows.meta,positioning,calc,shadows.blur,fit,backgrounds}
\usepackage{tcolorbox}
\tcbuselibrary{skins,breakable}
\usepackage{fontawesome5}
\usepackage{enumitem}
\usepackage{hyperref}
\hypersetup{colorlinks=true,linkcolor=indigo,urlcolor=indigo,pdfauthor={MonEcole},pdftitle={Manuel --- Module Enseignements}}
\usepackage{fancyhdr}
\usepackage{tabularx}
\usepackage{booktabs}
\usepackage{float}
\usepackage{caption}

% ====== COULEURS ======
\definecolor{indigo}{HTML}{4338CA}
\definecolor{violet}{HTML}{7C3AED}
\definecolor{amber}{HTML}{D97706}
\definecolor{emerald}{HTML}{059669}
\definecolor{rose}{HTML}{E11D48}
\definecolor{slate}{HTML}{334155}
\definecolor{slatelight}{HTML}{94A3B8}
\definecolor{teal}{HTML}{0D9488}
\definecolor{sky}{HTML}{0EA5E9}
\definecolor{coverbg}{HTML}{1E1B4B}
\definecolor{coveraccent}{HTML}{10B981}
\definecolor{lightbg}{HTML}{F8FAFC}
\definecolor{warningbg}{HTML}{FEF3C7}
\definecolor{successbg}{HTML}{ECFDF5}
\definecolor{infobg}{HTML}{EFF6FF}
\definecolor{dangerbg}{HTML}{FEF2F2}
\definecolor{purplebg}{HTML}{FAF5FF}
\definecolor{tealbg}{HTML}{F0FDFA}

% ====== TCOLORBOX ======
\newtcolorbox{astuce}[1][]{enhanced,colback=successbg,colframe=emerald,fonttitle=\bfseries\small,title={\faLightbulb\hspace{6pt}Astuce},left=8pt,right=8pt,top=6pt,bottom=6pt,arc=4pt,boxrule=1pt,breakable,#1}
\newtcolorbox{attention}[1][]{enhanced,colback=warningbg,colframe=amber,fonttitle=\bfseries\small,title={\faExclamationTriangle\hspace{6pt}Attention},left=8pt,right=8pt,top=6pt,bottom=6pt,arc=4pt,boxrule=1pt,breakable,#1}
\newtcolorbox{info}[1][]{enhanced,colback=infobg,colframe=indigo,fonttitle=\bfseries\small,title={\faInfoCircle\hspace{6pt}Information},left=8pt,right=8pt,top=6pt,bottom=6pt,arc=4pt,boxrule=1pt,breakable,#1}
\newtcolorbox{danger}[1][]{enhanced,colback=dangerbg,colframe=rose,fonttitle=\bfseries\small,title={\faTimes\hspace{6pt}Important},left=8pt,right=8pt,top=6pt,bottom=6pt,arc=4pt,boxrule=1pt,breakable,#1}
\newtcolorbox{etape}[1][]{enhanced,colback=white,colframe=indigo!60,fonttitle=\bfseries\small,left=10pt,right=10pt,top=8pt,bottom=8pt,arc=6pt,boxrule=1.5pt,breakable,#1}
\newtcolorbox{prerequis}[1][]{enhanced,colback=purplebg,colframe=violet,fonttitle=\bfseries\small,title={\faClipboard\hspace{6pt}Pr\'erequis},left=8pt,right=8pt,top=6pt,bottom=6pt,arc=4pt,boxrule=1pt,breakable,#1}
\newtcolorbox{resultat}[1][]{enhanced,colback=tealbg,colframe=teal,fonttitle=\bfseries\small,title={\faCheck\hspace{6pt}R\'esultat attendu},left=8pt,right=8pt,top=6pt,bottom=6pt,arc=4pt,boxrule=1pt,breakable,#1}

% ====== HEADERS ======
\pagestyle{fancy}\fancyhf{}
\fancyhead[L]{\small\color{slatelight}\leftmark}
\fancyhead[R]{\small\color{emerald}\textbf{MonEcole}}
\fancyfoot[C]{\small\color{slatelight}\thepage}
\renewcommand{\headrulewidth}{0.5pt}
\renewcommand{\footrulewidth}{0.3pt}

% ====== COMMANDES ======
\newcommand{\bouton}[1]{{\colorbox{emerald}{\small\sffamily\bfseries\color{white}\,#1\,}}}
\newcommand{\onglet}[1]{{\colorbox{indigo}{\small\sffamily\bfseries\color{white}\,#1\,}}}
\newcommand{\champ}[1]{\texttt{\color{violet}#1}}
\newcommand{\btnviolet}[1]{{\colorbox{violet}{\small\sffamily\bfseries\color{white}\,#1\,}}}
\newcommand{\btnamber}[1]{{\colorbox{amber}{\small\sffamily\bfseries\color{white}\,#1\,}}}
\newcommand{\btnrose}[1]{{\colorbox{rose}{\small\sffamily\bfseries\color{white}\,#1\,}}}

\begin{document}

% ===================================================================
%                         PAGE DE COUVERTURE
% ===================================================================
\begin{titlepage}
\begin{tikzpicture}[remember picture,overlay]
  \fill[coverbg] (current page.south west) rectangle (current page.north east);
  \fill[coveraccent,opacity=0.12] (current page.north east) ++(-4,1) circle (8cm);
  \fill[teal,opacity=0.08] (current page.south west) ++(3,-2) circle (10cm);
  \fill[coveraccent,opacity=0.06] (current page.center) ++(6,-4) circle (5cm);
  \fill[amber,opacity=0.05] (current page.north west) ++(2,-8) circle (4cm);
  \draw[coveraccent,line width=2pt,opacity=0.4] ([yshift=-6cm]current page.north west) ++(2.5,0) -- ++(16,0);
  \node[fill=coveraccent,rounded corners=14pt,minimum size=2cm,inner sep=0pt] at ([yshift=-3.5cm]current page.north) {\color{white}\fontsize{28pt}{28pt}\selectfont\faChalkboardTeacher};
  \node[anchor=north,text width=14cm,align=center] at ([yshift=-6.5cm]current page.north) {
    \color{white}\fontsize{32pt}{36pt}\selectfont\textbf{Manuel d'Utilisation}\\[12pt]
    \fontsize{18pt}{22pt}\selectfont\color{coveraccent!70!white}Module \textbf{Enseignements}
  };
  \node[anchor=north,text width=14cm,align=center] at ([yshift=-10cm]current page.north) {
    \color{white!60}\fontsize{11pt}{15pt}\selectfont
    Guide complet pour le tableau de bord des enseignants,\\
    la gestion des tables de r\'ef\'erence, l'attribution des cours\\
    aux enseignants et la planification des emplois du temps.
  };
  \draw[coveraccent,line width=1.5pt] ([yshift=-12.5cm]current page.north) ++(-3,0) -- ++(6,0);
  \node[anchor=north,text width=14cm,align=center] at ([yshift=-13.5cm]current page.north) {
    \color{white!50}\small
    \faBuilding\hspace{4pt}Plateforme MonEcole --- Architecture Hub-and-Spoke\\[6pt]
    \faCalendar\hspace{4pt}Avril 2026 --- Version 1.0\\[6pt]
    \faBook\hspace{4pt}Documentation Utilisateur --- 4 Sections D\'etaill\'ees
  };
  \node[anchor=south,text width=14cm,align=center] at ([yshift=1.5cm]current page.south) {
    \color{white!30}\footnotesize \faCopyright\hspace{3pt} Nexora Tech --- Tous droits r\'eserv\'es
  };
\end{tikzpicture}
\end{titlepage}

\tableofcontents
\newpage

% ===================================================================
\chapter{Introduction au Module Enseignements}
% ===================================================================

\section{Pr\'esentation g\'en\'erale}

Le module \textbf{Enseignements} est le centre de gestion du personnel enseignant et de l'organisation p\'edagogique de votre \'etablissement. Il couvre l'ensemble des op\'erations li\'ees aux enseignants~: de l'enregistrement dans la base de donn\'ees jusqu'\`a la planification de leur emploi du temps hebdomadaire.

Ce module est accessible depuis la barre de navigation lat\'erale (sidebar) en cliquant sur \textbf{Enseignements}. Il se divise en \textbf{4 sections} principales, chacune accessible via un lien dans la barre lat\'erale~:

\begin{center}
\begin{tikzpicture}[
  sect/.style={rounded corners=10pt, minimum height=1.1cm, minimum width=3.5cm, font=\small\bfseries, inner sep=8pt, text=white, drop shadow={shadow xshift=1pt, shadow yshift=-1pt, opacity=0.2}},
  desc/.style={text width=3.3cm, align=center, font=\tiny\color{slate}}
]
  \node[sect, fill=emerald] (s1) at (0,0) {\faChartBar\ Dashboard};
  \node[sect, fill=violet, right=6pt of s1] (s2) {\faDatabase\ R\'ef\'erence};
  \node[sect, fill=rose, right=6pt of s2] (s3) {\faUserTag\ Attribution};
  \node[sect, fill=amber, right=6pt of s3] (s4) {\faCalendar\ Horaire};

  \node[desc, below=8pt of s1] {Vue d'ensemble statistique des enseignants et du personnel};
  \node[desc, below=8pt of s2] {Tables de r\'ef\'erence~: cat\'egories, types, postes, dipl\^omes...};
  \node[desc, below=8pt of s3] {Attribuer un enseignant \`a chaque cours d'une classe};
  \node[desc, below=8pt of s4] {Planifier l'emploi du temps hebdomadaire par classe};
\end{tikzpicture}
\end{center}

\section{Flux de travail recommand\'e}

Pour un \'etablissement nouvellement configur\'e, voici l'ordre logique des op\'erations~:

\begin{center}
\begin{tikzpicture}[
  step/.style={draw=slate!40, fill=white, rounded corners=8pt, minimum width=2.4cm, minimum height=1.5cm, align=center, font=\scriptsize\bfseries, text=slate, drop shadow={shadow xshift=1pt, shadow yshift=-1pt, opacity=0.12}},
  arrow/.style={-{Stealth[length=6pt,width=5pt]}, thick, color=slatelight}
]
  \node[step] (s1) at (0,0) {\faDatabase\\[2pt]Configurer les\\tables de r\'ef.};
  \node[step, right=0.5cm of s1] (s2) {\faUserPlus\\[2pt]Enregistrer les\\enseignants};
  \node[step, right=0.5cm of s2] (s3) {\faUserTag\\[2pt]Attribuer les\\cours};
  \node[step, right=0.5cm of s3] (s4) {\faCalendar\\[2pt]Planifier\\l'horaire};
  \node[step, right=0.5cm of s4] (s5) {\faChartBar\\[2pt]Suivre via\\le dashboard};

  \draw[arrow] (s1) -- (s2);
  \draw[arrow] (s2) -- (s3);
  \draw[arrow] (s3) -- (s4);
  \draw[arrow] (s4) -- (s5);
\end{tikzpicture}
\end{center}

\begin{info}
Les sections \textbf{R\'ef\'erence} et \textbf{Dashboard} sont des modules de consultation/configuration de base. L'\textbf{Attribution} et l'\textbf{Horaire} sont les modules op\'erationnels principaux o\`u se d\'eroule le travail quotidien de planification.
\end{info}

\newpage
% ===================================================================
\chapter{Section 2 --- R\'ef\'erence du Personnel}
% ===================================================================

\section{Objectif de cette section}

La section \textbf{R\'ef\'erence} permet de g\'erer les \textbf{tables de r\'ef\'erence} utilis\'ees pour cat\'egoriser et d\'ecrire le personnel enseignant. Ces tables sont essentielles car elles alimentent les menus d\'eroulants utilis\'es lors de l'enregistrement des enseignants.

\section{Les 7 tables de r\'ef\'erence}

La section affiche une barre de boutons color\'es, chacun correspondant \`a une table de r\'ef\'erence~:

\begin{center}
\begin{tikzpicture}[
  btn/.style={rounded corners=4pt, minimum height=0.6cm, font=\tiny\bfseries, inner sep=4pt, text=white, fill=#1}
]
  \node[btn=indigo] (b1) at (0,0) {\faTags\ Cat\'egories};
  \node[btn=violet, right=3pt of b1] (b2) {\faUserTag\ Types};
  \node[btn=teal, right=3pt of b2] (b3) {\faBriefcase\ Postes};
  \node[btn=sky, right=3pt of b3] (b4) {\faGraduationCap\ Dipl\^omes};
  \node[btn=amber, right=3pt of b4] (b5) {\faMicroscope\ Sp\'ecialit\'es};
  \node[btn=rose, right=3pt of b5] (b6) {\faClock\ Vacations};
  \node[btn=slate, right=3pt of b6] (b7) {\faTasks\ T\^aches};
\end{tikzpicture}
\end{center}

\begin{center}
\begin{tabularx}{\textwidth}{llX}
\toprule
\textbf{Table} & \textbf{Ic\^one} & \textbf{Description} \\
\midrule
Cat\'egories & \faTags & Classement professionnel~: Enseignant, Administratif, Personnel d'appui... \\
Types & \faUserTag & Sous-classification~: Titulaire, Vacataire, Stagiaire, Contractuel... \\
Postes & \faBriefcase & Fonctions~: Professeur, Directeur, Pr\'efet d'\'etudes, Surveillant... \\
Dipl\^omes & \faGraduationCap & Qualifications~: Licence, Master, Doctorat, Brevet, Certificat... \\
Sp\'ecialit\'es & \faMicroscope & Domaines d'expertise~: Math\'ematiques, Fran\c{c}ais, Physique, Biologie... \\
Vacations & \faClock & Types de vacation~: Temps plein, Mi-temps, Horaire... \\
T\^aches & \faTasks & Activit\'es secondaires~: Titulaire de classe, Chef de d\'epartement... \\
\bottomrule
\end{tabularx}
\end{center}

\section{Pour consulter une table de r\'ef\'erence}

\begin{etape}[title=Proc\'edure : Consulter une table]
\begin{enumerate}[leftmargin=1.5cm,label=\textbf{\arabic*.},itemsep=8pt]
  \item \textbf{Acc\'edez \`a la section R\'ef\'erence} en cliquant sur \textbf{R\'ef\'erence} dans la barre lat\'erale.
  \item \textbf{Cliquez sur le bouton} correspondant \`a la table souhait\'ee (ex.\ : \onglet{Cat\'egories}). Le bouton s\'electionn\'e prend une couleur indigo.
  \item \textbf{Le contenu s'affiche} dans un tableau structur\'e en dessous des boutons, avec les colonnes \textbf{ID}, \textbf{Nom}, \textbf{Description} (si applicable) et \textbf{Actions} (modifier / supprimer).
\end{enumerate}
\end{etape}

\section{Pour ajouter un \'el\'ement \`a une table}

\begin{etape}[title=Proc\'edure : Ajouter un \'el\'ement]
\begin{enumerate}[leftmargin=1.5cm,label=\textbf{\arabic*.},itemsep=10pt]
  \item \textbf{S\'electionnez la table} via les boutons (ex.\ : \onglet{Postes}).
  \item \textbf{Cliquez sur} \bouton{+ Ajouter} en haut du tableau. Un formulaire inline appara\^it.
  \item \textbf{Remplissez le nom} de l'\'el\'ement (ex.\ : \textit{Professeur principal}).
  \item \textbf{Ajoutez une description} (optionnel) pour pr\'eciser le r\^ole.
  \item \textbf{Cliquez sur} \bouton{Enregistrer}. L\'el\'ement appara\^it dans le tableau.
\end{enumerate}
\end{etape}

\section{Pour modifier ou supprimer un \'el\'ement}

\begin{etape}[title=Proc\'edure : Modifier / Supprimer]
\begin{enumerate}[leftmargin=1.5cm,label=\textbf{\arabic*.},itemsep=8pt]
  \item \textbf{Rep\'erez la ligne} de l'\'el\'ement dans le tableau.
  \item \textbf{Pour modifier}~: cliquez sur \textcolor{indigo}{\faPen}. Les champs deviennent \'editables. Modifiez et cliquez sur \bouton{OK}.
  \item \textbf{Pour supprimer}~: cliquez sur \textcolor{rose}{\faTrash}. Confirmez dans la bo\^ite de dialogue.
\end{enumerate}
\end{etape}

\begin{attention}
La suppression d'un \'el\'ement de r\'ef\'erence peut impacter les enseignants qui y sont rattach\'es. Par exemple, supprimer la cat\'egorie \textit{<<~Vacataire~>>} d\'etachera tous les enseignants class\'es dans cette cat\'egorie.
\end{attention}

\section{Gestion compl\`ete du personnel}

La section R\'ef\'erence contient \'egalement le \textbf{panneau de gestion du personnel}, accessible via un bouton dans la barre d'outils~:

\subsection{Pour ajouter un enseignant}

\begin{etape}[title=Proc\'edure d\'etaill\'ee : Ajouter un enseignant]
\begin{enumerate}[leftmargin=1.5cm,label=\textbf{\arabic*.},itemsep=10pt]
  \item \textbf{Cliquez sur} \bouton{Ajouter} (bouton vert avec ic\^one \faUserPlus).
  \item \textbf{Le formulaire d'ajout} s'affiche avec 6 champs sur 3 colonnes~:
  \begin{itemize}
    \item \champ{Nom} --- Nom de famille (obligatoire)
    \item \champ{Pr\'enom} --- Pr\'enom (obligatoire)
    \item \champ{Post-nom} --- Post-nom (optionnel)
    \item \champ{Genre} --- Menu d\'eroulant~: Masculin / F\'eminin (obligatoire)
    \item \champ{T\'el\'ephone} --- Num\'ero de t\'el\'ephone
    \item \champ{Email} --- Adresse \'electronique
  \end{itemize}
  \item \textbf{Cliquez sur} \bouton{Enregistrer}. Un toast vert confirme l'ajout et la grille se rafra\^ichit.
\end{enumerate}
\end{etape}

\subsection{Pour modifier un enseignant dans la grille}

\begin{etape}[title=Proc\'edure : Modification inline]
\begin{enumerate}[leftmargin=1.5cm,label=\textbf{\arabic*.},itemsep=8pt]
  \item Rep\'erez l'enseignant dans la grille du personnel.
  \item \textbf{Cliquez directement sur la cellule} que vous souhaitez modifier (les cellules sont \'editables en mode \texttt{contenteditable}).
  \item La cellule prend un contour bleu indigo pour signaler le mode \'edition.
  \item \textbf{Tapez votre modification} puis cliquez en dehors de la cellule ou appuyez sur Tab.
  \item La cellule \textbf{clignote en vert} bri\`evement pour confirmer l'enregistrement automatique.
\end{enumerate}
\end{etape}

\subsection{Pour importer du personnel depuis Excel}

\begin{etape}[title=Proc\'edure : Import Excel]
\begin{enumerate}[leftmargin=1.5cm,label=\textbf{\arabic*.},itemsep=10pt]
  \item \textbf{T\'el\'echargez le mod\`ele} en cliquant sur \btnviolet{Mod\`ele Import} (bouton violet).
  \item \textbf{Remplissez le fichier} Excel t\'el\'echarg\'e avec les donn\'ees du personnel~: nom, pr\'enom, post-nom, genre, t\'el\'ephone, email.
  \item \textbf{Cliquez sur} \btnrose{Importer} (bouton rose) et s\'electionnez le fichier rempli.
  \item Le syst\`eme importe les donn\'ees et affiche un r\'esum\'e~: nombre de nouveaux enregistrements, nombre de doublons ignor\'es.
\end{enumerate}
\end{etape}

\subsection{Pour personnaliser les colonnes affich\'ees}

\begin{etape}[title=Proc\'edure : S\'electeur de colonnes]
\begin{enumerate}[leftmargin=1.5cm,label=\textbf{\arabic*.},itemsep=8pt]
  \item Cliquez sur le bouton \onglet{Colonnes} (en haut \`a droite de la grille).
  \item Un panneau flottant appara\^it avec la liste de toutes les colonnes disponibles, chacune avec une case \`a cocher.
  \item \textbf{Cochez ou d\'ecochez} les colonnes pour les afficher ou les masquer.
  \item La grille se met \`a jour imm\'ediatement.
\end{enumerate}
\end{etape}

\newpage
% ===================================================================
\chapter{Section 3 --- Attribution des Cours}
% ===================================================================

\section{Objectif de cette section}

La section \textbf{Attribution} est le c\oe ur op\'erationnel du module Enseignements. Elle permet d'\textbf{associer un enseignant \`a chaque cours} pour chaque classe de l'\'etablissement. Sans attribution, le syst\`eme ne peut pas g\'en\'erer les emplois du temps ni associer les notes aux enseignants.

\begin{info}
L'attribution se fait \textbf{classe par classe}. Pour chaque classe, vous verrez la liste de tous les cours configur\'es, et vous pourrez assigner un enseignant \`a chacun d'eux \`a l'aide d'un menu d\'eroulant.
\end{info}

\section{Comprendre l'interface}

L'interface se compose de~:

\begin{enumerate}[leftmargin=1.2cm]
  \item \textbf{S\'electeur de classe} --- Un menu d\'eroulant dans un encadr\'e violet avec ic\^one \faChalkboard. La s\'election d'une classe charge la liste des cours.
  \item \textbf{Compteurs r\'ecapitulatifs} --- 3 cartes affichant~:
  \begin{itemize}
    \item \textcolor{emerald}{\textbf{Total cours}} --- Nombre total de cours de la classe.
    \item \textcolor{sky}{\textbf{Attribu\'es}} --- Nombre de cours d\'ej\`a affect\'es \`a un enseignant.
    \item \textcolor{rose}{\textbf{Non attribu\'es}} --- Nombre de cours sans enseignant.
  \end{itemize}
  \item \textbf{Tableau d'attribution} --- Liste de tous les cours avec colonnes~: Cours, Domaine, Heures, Enseignant (menu d\'eroulant), Actions.
\end{enumerate}



\section{Pour attribuer un enseignant \`a un cours}

\begin{prerequis}
\begin{itemize}
  \item La \textbf{configuration annuelle} doit \^etre compl\`ete (cycles, classes, cours activ\'es).
  \item Des \textbf{enseignants} doivent avoir \'et\'e enregistr\'es dans la section R\'ef\'erence.
\end{itemize}
\end{prerequis}

\begin{etape}[title=Proc\'edure d\'etaill\'ee : Attribuer un enseignant]
\begin{enumerate}[leftmargin=1.5cm,label=\textbf{\arabic*.},itemsep=12pt]
  \item \textbf{S\'electionnez une classe} dans le menu d\'eroulant \champ{Classe}. Le tableau des cours se charge avec les compteurs r\'ecapitulatifs.

  \item \textbf{Rep\'erez le cours} \`a attribuer dans le tableau. Les cours non attribu\'es sont marqu\'es d'un indicateur rouge.

  \item \textbf{Ouvrez le menu d\'eroulant} dans la colonne \textbf{Enseignant} de la ligne correspondante. La liste affiche tous les enseignants enregistr\'es, tri\'es alphab\'etiquement.

  \item \textbf{S\'electionnez l'enseignant} souhait\'e. Le syst\`eme enregistre imm\'ediatement l'attribution via l'API.

  \item \textbf{Le compteur se met \`a jour} automatiquement~: le nombre d'\textit{<<~Attribu\'es~>>} augmente et celui de \textit{<<~Non attribu\'es~>>} diminue.

  \item \textbf{Pour retirer une attribution}~: s\'electionnez l'option vide (premi\`ere option du menu d\'eroulant) ou cliquez sur \textcolor{rose}{\faTimes} \`a c\^ot\'e du nom.
\end{enumerate}
\end{etape}

\begin{resultat}
Apr\`es attribution, l'enseignant est visible dans le tableau avec son nom complet. L'information est imm\'ediatement disponible pour~:
\begin{itemize}
  \item La planification des emplois du temps (section Horaire).
  \item Les grilles de notes du module \'Evaluations.
  \item La g\'en\'eration des bulletins PDF.
\end{itemize}
\end{resultat}

\begin{astuce}
\textbf{Un m\^eme enseignant} peut \^etre attribu\'e \`a plusieurs cours et \`a plusieurs classes. Le syst\`eme v\'erifiera automatiquement les conflits d'horaire lors de la planification de l'emploi du temps.
\end{astuce}

\section{Pour visualiser les attributions existantes}

\begin{etape}[title=Proc\'edure : V\'erifier les attributions]
\begin{enumerate}[leftmargin=1.5cm,label=\textbf{\arabic*.},itemsep=8pt]
  \item S\'electionnez une classe dans le menu d\'eroulant.
  \item Le tableau se charge. Les cours d\'ej\`a attribu\'es affichent le nom de l'enseignant dans la colonne \textbf{Enseignant} avec un fond vert clair.
  \item Les cours non attribu\'es affichent \textit{<<~S\'electionnez...~>>} en gris dans le menu d\'eroulant.
  \item Utilisez les 3 compteurs pour \'evaluer rapidement la progression.
\end{enumerate}
\end{etape}

\newpage
% ===================================================================
\chapter{Section 4 --- Emploi du Temps (Horaire)}
% ===================================================================

\section{Objectif de cette section}

La section \textbf{Horaire} permet de planifier l'\textbf{emploi du temps hebdomadaire} de chaque classe. Elle offre une vue calendaire interactive o\`u vous pouvez cr\'eer, modifier et supprimer des cr\'eneaux de cours pour chaque jour de la semaine.

\begin{info}
L'emploi du temps est organis\'e \textbf{par semaine}. Vous pouvez naviguer entre les semaines avec les boutons de navigation (\faChevronLeft\ et \faChevronRight) et copier l'emploi du temps d'une semaine vers la suivante.
\end{info}

\section{Comprendre l'interface}

L'interface se compose de~:

\begin{enumerate}[leftmargin=1.2cm]
  \item \textbf{S\'electeur de classe} --- Menu d\'eroulant avec fond violet.
  \item \textbf{S\'electeur de p\'eriode} --- Menu d\'eroulant avec fond jaune/ambre. D\'efinit la plage de dates pour les semaines disponibles.
  \item \textbf{Navigation entre les semaines} --- Boutons \faChevronLeft\ / \faChevronRight\ avec le libell\'e de la semaine (ex.\ : \textit{<<~Semaine du 14 au 18 octobre 2025~>>}).
  \item \textbf{Bouton Copier semaine pr\'ec\'edente} --- Permet de dupliquer tous les cr\'eneaux de la semaine pr\'ec\'edente.
  \item \textbf{Grille hebdomadaire} --- Un tableau avec les jours en colonnes et les heures en lignes. Chaque cellule repr\'esente un cr\'eneau potentiel.
\end{enumerate}

\section{Pour afficher l'emploi du temps d'une classe}

\begin{etape}[title=Proc\'edure : Charger l'emploi du temps]
\begin{enumerate}[leftmargin=1.5cm,label=\textbf{\arabic*.},itemsep=12pt]
  \item \textbf{S\'electionnez une classe} dans le premier menu d\'eroulant \champ{Classe}.

  \item \textbf{S\'electionnez une p\'eriode} dans le second menu d\'eroulant \champ{P\'eriode}. Le syst\`eme charge les semaines disponibles pour cette p\'eriode.

  \item \textbf{La navigation par semaine} appara\^it avec le libell\'e de la premi\`ere semaine et un compteur (ex.\ : \textit{<<~Semaine 1/8~>>}).

  \item \textbf{La grille hebdomadaire} se charge avec les cr\'eneaux existants. Chaque cr\'eneau est repr\'esent\'e par un bloc color\'e affichant le nom du cours, l'heure de d\'ebut et de fin, et le nom de l'enseignant attribu\'e.
\end{enumerate}
\end{etape}

\section{Pour ajouter un cr\'eneau de cours}

\begin{etape}[title=Proc\'edure d\'etaill\'ee : Ajouter un cr\'eneau]
\begin{enumerate}[leftmargin=1.5cm,label=\textbf{\arabic*.},itemsep=12pt]
  \item \textbf{Cliquez sur une cellule vide} de la grille (correspondant au jour et \`a l'heure souhait\'es). La fen\^etre modale \textit{<<~Ajouter un cr\'eneau~>>} s'ouvre.

  \item \textbf{La date est pr\'e-remplie} avec le jour s\'electionn\'e. Elle est affich\'ee en lecture seule dans le champ \champ{Date}.

  \item \textbf{S\'electionnez le cours} dans le menu d\'eroulant \champ{Cours}. Seuls les cours de la classe s\'electionn\'ee sont list\'es (avec le nom de l'enseignant attribu\'e s'il existe).

  \item \textbf{D\'efinissez les heures}~:
  \begin{itemize}
    \item \champ{D\'ebut} --- S\'electionnez l'heure de d\'ebut (format HH:MM).
    \item \champ{Fin} --- S\'electionnez l'heure de fin.
  \end{itemize}

  \item \textbf{Option multi-jours} (optionnel)~: Cochez la case \textit{<<~Appliquer \`a plusieurs jours~>>} pour cr\'eer le m\^eme cr\'eneau sur d'autres jours de la semaine. Les jours disponibles s'affichent sous forme de boutons cliquables (Lundi, Mardi, etc.). S\'electionnez les jours souhait\'es.

  \item \textbf{Cliquez sur} \bouton{Enregistrer}. Le cr\'eneau appara\^it dans la grille.
\end{enumerate}
\end{etape}

\begin{astuce}
L'option \textbf{multi-jours} est particuli\`erement utile pour les cours r\'ecurrents. Par exemple, si les Math\'ematiques ont lieu tous les lundis et mercredis de 8h \`a 10h, vous n'avez besoin de cr\'eer qu'un seul cr\'eneau en s\'electionnant les deux jours.
\end{astuce}

\section{Pour modifier un cr\'eneau existant}

\begin{etape}[title=Proc\'edure : Modifier un cr\'eneau]
\begin{enumerate}[leftmargin=1.5cm,label=\textbf{\arabic*.},itemsep=8pt]
  \item \textbf{Cliquez sur le bloc color\'e} du cr\'eneau dans la grille.
  \item La modale s'ouvre \textbf{pr\'e-remplie} avec les informations existantes (cours, heures).
  \item Modifiez les champs n\'ecessaires.
  \item Cliquez sur \bouton{Enregistrer}.
\end{enumerate}
\end{etape}

\section{Pour supprimer un cr\'eneau}

\begin{etape}[title=Proc\'edure : Supprimer un cr\'eneau]
\begin{enumerate}[leftmargin=1.5cm,label=\textbf{\arabic*.},itemsep=8pt]
  \item Cliquez sur le cr\'eneau dans la grille.
  \item Dans la modale, cliquez sur \btnrose{Supprimer} (bouton rouge en bas).
  \item Confirmez la suppression. Le cr\'eneau dispara\^it de la grille.
\end{enumerate}
\end{etape}

\section{Pour copier l'emploi du temps d'une semaine pr\'ec\'edente}

\begin{etape}[title=Proc\'edure : Copier une semaine]
\begin{enumerate}[leftmargin=1.5cm,label=\textbf{\arabic*.},itemsep=10pt]
  \item \textbf{Naviguez vers la semaine cible} (celle qui doit recevoir les cr\'eneaux) \`a l'aide des boutons \faChevronLeft\ / \faChevronRight.

  \item \textbf{Cliquez sur} \bouton{Copier semaine pr\'ec\'edente} (bouton vert avec ic\^one \faCopy, visible uniquement \`a partir de la semaine 2).

  \item \textbf{La modale de copie s'ouvre} et affiche la liste de tous les cr\'eneaux de la semaine pr\'ec\'edente, chacun avec une case \`a cocher. Vous pouvez~:
  \begin{itemize}
    \item \textbf{Tout s\'electionner} --- Cliquez sur \bouton{Tout s\'electionner} pour cocher tous les cr\'eneaux.
    \item \textbf{S\'electionner individuellement} --- Cochez uniquement les cr\'eneaux que vous souhaitez copier.
    \item \textbf{Tout d\'es\'electionner} --- Pour recommencer la s\'election.
  \end{itemize}

  \item \textbf{Cliquez sur} \bouton{Copier la s\'election}. Les cr\'eneaux s\'electionn\'es sont dupliqu\'es sur les m\^emes jours/heures de la semaine courante.
\end{enumerate}
\end{etape}

\begin{danger}
La copie de semaine \textbf{n'\'ecrase pas} les cr\'eneaux existants. Si un cr\'eneau identique existe d\'ej\`a, le syst\`eme l'ignore. Les cr\'eneaux sont cr\'e\'es avec les m\^emes heures mais aux dates correspondantes de la nouvelle semaine.
\end{danger}

\newpage
% ===================================================================
\chapter{R\'ef\'erence rapide}
% ===================================================================

\section{Navigation entre les sections}

\begin{tabularx}{\textwidth}{llX}
\toprule
\textbf{Ic\^one Sidebar} & \textbf{Section} & \textbf{Description} \\
\midrule
\faChartBar\ (vert) & Dashboard & Tableau de bord statistique du personnel \\
\faDatabase\ (violet) & R\'ef\'erence & Tables de r\'ef\'erence + gestion du personnel \\
\faUserTag\ (rouge) & Attribution & Attribution enseignant $\rightarrow$ cours par classe \\
\faCalendar\ (ambre) & Horaire & Planification de l'emploi du temps hebdomadaire \\
\bottomrule
\end{tabularx}

\section{Endpoints API du module Enseignements}

{\small
\begin{tabularx}{\textwidth}{lX}
\toprule
\textbf{Endpoint} & \textbf{Description} \\
\midrule
\texttt{/api/personnel/list/} & GET --- Liste compl\`ete du personnel \\
\texttt{/api/personnel/add/} & POST --- Ajouter un membre du personnel \\
\texttt{/api/personnel/update/} & POST --- Modifier les donn\'ees d'un personnel \\
\texttt{/api/personnel/delete/} & POST --- Supprimer un personnel \\
\texttt{/api/personnel/import/} & POST --- Importer depuis Excel \\
\texttt{/api/personnel/template/} & GET --- T\'el\'echarger le mod\`ele Excel \\
\texttt{/api/ref-tables/get/} & GET --- Lister les \'el\'ements d'une table \\
\texttt{/api/ref-tables/save/} & POST --- Ajouter/modifier un \'el\'ement \\
\texttt{/api/ref-tables/delete/} & POST --- Supprimer un \'el\'ement de r\'ef\'erence \\
\texttt{/api/attribution/list/} & GET --- Liste des attributions d'une classe \\
\texttt{/api/attribution/save/} & POST --- Attribuer un enseignant \`a un cours \\
\texttt{/api/attribution/remove/} & POST --- Retirer une attribution \\
\texttt{/api/horaire/get/} & GET --- Cr\'eneaux d'une classe/semaine \\
\texttt{/api/horaire/save/} & POST --- Cr\'eer/modifier un cr\'eneau \\
\texttt{/api/horaire/delete/} & POST --- Supprimer un cr\'eneau \\
\texttt{/api/horaire/copy-week/} & POST --- Copier les cr\'eneaux d'une semaine \\
\texttt{/api/dashboard/eleves-stats/} & GET --- Statistiques \'el\`eves \\
\texttt{/api/dashboard/campus-list/} & GET --- Liste des campus \\
\bottomrule
\end{tabularx}
}

\section{Glossaire}

\begin{description}[leftmargin=1.5cm, style=nextline, font=\bfseries\color{indigo}]
  \item[Attribution] Lien entre un enseignant et un cours pour une classe donn\'ee dans l'ann\'ee scolaire.
  \item[Campus] Site physique de l'\'etablissement (un \'etablissement peut avoir plusieurs campus).
  \item[Cat\'egorie] Classification professionnelle du personnel (Enseignant, Administratif...).
  \item[Cr\'eneau] Bloc horaire d'un cours planifi\'e \`a un jour et une heure pr\'ecise.
  \item[Emploi du temps] Grille hebdomadaire des cr\'eneaux de cours pour une classe.
  \item[P\'eriode] Subdivision du calendrier scolaire d\'elimitant la plage de l'emploi du temps.
  \item[Table de r\'ef\'erence] Liste de valeurs pr\'ed\'efinies utilis\'ees pour classifier le personnel.
  \item[Vacation] R\'egime de travail d'un enseignant (temps plein, mi-temps, horaire).
\end{description}

\vspace{2cm}
\begin{center}
\begin{tikzpicture}
  \fill[coverbg, rounded corners=12pt] (0,0) rectangle (14,2.5);
  \node[text=white, font=\large\bfseries] at (7,1.5) {\faCheck\hspace{8pt}Fin du Manuel --- Module Enseignements};
  \node[text=white!50, font=\small] at (7,0.7) {Plateforme MonEcole --- Nexora Tech --- Avril 2026 --- v1.0};
\end{tikzpicture}
\end{center}

\end{document}
